% =====================================================================
\chapter{Conjuntos y operaciones}
\label{tema:1}
% =====================================================================

Este tema fija el vocabulario básico de la teoría de conjuntos y, sobre todo, la técnica para demostrar que dos conjuntos son iguales. La presentación sigue el capítulo~1 y el capítulo~8 de Hammack~\cite{hammack}, las secciones~1.3, 1.4 y~2.3 de Velleman~\cite{velleman} y la sección~4.1 de~\cite{mcs}; el estilo de las demostraciones es el de la clase «Set Theory Revisited» de Stanford CS103~\cite{cs103set}.

\section{Conjuntos y pertenencia}

\begin{definicion}[Conjunto]
Un \emph{conjunto} es una colección de objetos, llamados \emph{elementos}, considerada como un objeto en sí misma. Se escribe $x\in A$ si $x$ es elemento de $A$ y $x\notin A$ en caso contrario. Dos conjuntos son \emph{iguales} si tienen exactamente los mismos elementos:
\[
A=B\iff \forall x\,(x\in A\leftrightarrow x\in B).
\]
\end{definicion}

Un conjunto puede describirse \emph{por extensión}, enumerando sus elementos, como $\{1,2,3\}$, o \emph{por comprensión}, mediante una propiedad, como $\{x\in\Z : x^{2}<5\}=\{-2,-1,0,1,2\}$. Por la definición de igualdad, el orden y las repeticiones no importan: $\{1,2,2\}=\{2,1\}$. El \emph{conjunto vacío} $\emptyset$ no tiene elementos. Si $A$ tiene un número finito de elementos, ese número es su \emph{cardinal}, $\abs{A}$.

\addtocounter{definicion}{1}%

\section{Subconjuntos y conjunto de partes}

\begin{definicion}[Subconjunto]
$A$ es \emph{subconjunto} de $B$, y se escribe $A\subseteq B$, si todo elemento de $A$ es elemento de $B$: $\forall x\,(x\in A\rightarrow x\in B)$. Si además $A\neq B$, se dice que $A$ es un subconjunto \emph{propio} y se escribe $A\subsetneq B$.
\end{definicion}

\begin{proposicion}\label{prop:t1-subconj}
Para conjuntos cualesquiera $A$, $B$ y $C$:
\begin{enumerate}
  \item $\emptyset\subseteq A$ y $A\subseteq A$;
  \item si $A\subseteq B$ y $B\subseteq C$, entonces $A\subseteq C$;
  \item $A=B$ si y solo si $A\subseteq B$ y $B\subseteq A$.
\end{enumerate}
\end{proposicion}

El apartado~(iii) es el \emph{método de la doble inclusión}: para probar $A=B$ se toma un elemento arbitrario de $A$ y se demuestra que está en $B$, y después al revés~\cite[§8.3]{hammack}.

\begin{definicion}[Conjunto de partes]
El \emph{conjunto de las partes} de $A$ es el conjunto de todos sus subconjuntos: $\Pa(A)=\{X : X\subseteq A\}$.
\end{definicion}

\begin{teorema}\label{teo:t1-partes}
Si $A$ es finito con $\abs{A}=n$, entonces $\abs{\Pa(A)}=2^{n}$.
\end{teorema}

\addtocounter{definicion}{1}%

\section{Operaciones con conjuntos}

\begin{definicion}[Operaciones]
Sean $A$ y $B$ subconjuntos de un conjunto universal $U$. Se definen
\begin{align*}
A\cup B&=\{x : x\in A\ \text{o}\ x\in B\} &&\text{(unión)},\\
A\cap B&=\{x : x\in A\ \text{y}\ x\in B\} &&\text{(intersección)},\\
A\setminus B&=\{x : x\in A\ \text{y}\ x\notin B\} &&\text{(diferencia)},\\
\overline{A}&=U\setminus A &&\text{(complemento)},\\
A\dsim B&=(A\setminus B)\cup(B\setminus A) &&\text{(diferencia simétrica)}.
\end{align*}
$A$ y $B$ son \emph{disjuntos} si $A\cap B=\emptyset$.
\end{definicion}

Cada operación corresponde a un conector lógico: la unión a la disyunción, la intersección a la conjunción y el complemento a la negación. Por eso las leyes de la lógica proposicional se traducen en leyes de conjuntos~\cite[§1.4]{velleman}.

\begin{teorema}[Leyes del álgebra de conjuntos]\label{teo:t1-leyes}
Para subconjuntos $A$, $B$, $C$ de $U$:
\begin{enumerate}
  \item conmutativas: $A\cup B=B\cup A$, \ $A\cap B=B\cap A$;
  \item asociativas: $(A\cup B)\cup C=A\cup(B\cup C)$, \ $(A\cap B)\cap C=A\cap(B\cap C)$;
  \item distributivas: $A\cap(B\cup C)=(A\cap B)\cup(A\cap C)$, \ $A\cup(B\cap C)=(A\cup B)\cap(A\cup C)$;
  \item leyes de De Morgan: $\overline{A\cup B}=\overline{A}\cap\overline{B}$, \ $\overline{A\cap B}=\overline{A}\cup\overline{B}$;
  \item $A\setminus B=A\cap\overline{B}$ \ y \ $\overline{\overline{A}}=A$.
\end{enumerate}
\end{teorema}

\addtocounter{figure}{1}%

\addtocounter{definicion}{1}%

\addtocounter{definicion}{1}%

\section{Producto cartesiano y familias indexadas}

\begin{definicion}[Producto cartesiano]
Un \emph{par ordenado} $(a,b)$ es un objeto con primera componente $a$ y segunda $b$, con la propiedad $(a,b)=(c,d)\iff a=c\ \text{y}\ b=d$. El \emph{producto cartesiano} de $A$ y $B$ es $A\times B=\{(a,b): a\in A,\ b\in B\}$. Análogamente, $A_1\times\dots\times A_n$ es el conjunto de $n$-tuplas, y $A^{n}=A\times\dots\times A$ ($n$ factores).
\end{definicion}

\begin{proposicion}\label{prop:t1-producto}
Si $A$ y $B$ son finitos, $\abs{A\times B}=\abs{A}\cdot\abs{B}$; en particular, $\abs{A^{n}}=\abs{A}^{n}$. Además, para conjuntos cualesquiera, $A\times(B\cup C)=(A\times B)\cup(A\times C)$.
\end{proposicion}

\addtocounter{definicion}{1}%

\begin{definicion}[Familias indexadas]
Sea $I$ un conjunto de índices y, para cada $i\in I$, un conjunto $A_i$. Se definen
\[
\bigcup_{i\in I}A_i=\{x : x\in A_i\ \text{para algún}\ i\in I\},\qquad
\bigcap_{i\in I}A_i=\{x : x\in A_i\ \text{para todo}\ i\in I\}.
\]
\end{definicion}

\addtocounter{definicion}{1}%

\addtocounter{definicion}{1}%

